% ============================================================
% Section 04 — Identity and Non-Identity Through Difference
% Book I: Ontology of Viability
% Chapter 05: Identity and Continuity
%
% Purpose:
%   To examine whether difference between successive
%   conditions can establish numerical non-identity,
%   and to distinguish qualitative difference from
%   ontological discontinuity.
%
% Conceptual Policy:
%   1. Difference is not non-identity.
%   2. Qualitative transformation does not automatically
%      establish numerical replacement.
%   3. Similarity does not independently establish identity.
%   4. Difference must be specified by subject, domain,
%      and respect.
%   5. Neither identity nor non-identity should be inferred
%      from a single unqualified comparison.
% ============================================================

\section{Identity and Non-Identity Through Difference}
\label{sec:identity-and-non-identity-through-difference}

\subsection{The Ontological Problem of Difference and Non-Identity}

The relation between difference and non-identity constitutes
a fundamental problem in the investigation of identity
through transformation.

An existent may undergo qualitative, material, structural,
functional, relational, or manifestational differences
across successive conditions. Such differences may be
substantial and may alter the way in which the existent
is determined or manifested.

Nevertheless, the presence of difference does not
immediately establish numerical non-identity.

The proposition that an existent differs from a previous
condition and the proposition that it is no longer
the same existent are distinct ontological claims.
The first concerns a difference between determinations
or conditions. The second concerns the numerical identity
of the existent across the relevant succession.

Confusing these claims would make every transformation
a possible termination of identity merely because
some feature has changed. Such a conclusion would
require an independent argument establishing that
the relevant difference is incompatible with the
persistence of identity.

The opposite inference is equally problematic.
The absence of an identified difference, or the
presence of extensive similarity, does not independently
establish numerical identity.

The ontological problem is therefore not whether
difference exists, but whether a particular difference
has the significance required to establish non-identity.

\subsection{Qualitative Difference and Numerical Non-Identity}

Qualitative difference concerns the distinction between
features, determinations, properties, or conditions.

Numerical non-identity concerns whether two subjects
are distinct existents rather than one and the same
existent considered across different conditions.

These notions belong to different analytical levels.

Two conditions of an existent may differ qualitatively
while remaining conditions of one and the same existent.
A material configuration may change, a structural
organization may be modified, or a functional capacity
may be transformed without these changes alone
establishing numerical replacement.

Likewise, two numerically distinct existents may share
many or even all of the features selected for a particular
comparison.

Qualitative similarity and difference therefore cannot
be treated as direct measures of numerical identity.

The distinction requires that every comparison specify
whether it concerns features, configurations, relations,
manifestations, or the numerical identity of the
subjects being compared.

Without this distinction, qualitative descriptions
may be incorrectly transformed into ontological
conclusions.

\subsection{Difference Across Temporal Conditions}

The problem becomes particularly significant when
difference is examined across temporal conditions.

An existent may be described at different temporal
stages through distinct configurations or determinations.
The earlier and later conditions may differ in numerous
respects while still raising the question of whether
they belong to one and the same identity.

The mere fact that the conditions occur at different
times does not establish that they belong to different
existents. Nor does temporal succession establish
that they belong to one numerically identical existent.

The relation between successive conditions must be
examined rather than assumed.

A temporal sequence may involve transformation within
one existent, the succession of distinct existents,
or a more complex relation in which the relevant
subjects and respects have not yet been adequately
determined.

Consequently, the comparison of temporal stages
cannot be resolved through difference alone.

The temporal distance between two conditions, their
degree of similarity, or the magnitude of their
qualitative transformation does not independently
settle the question of numerical identity.

\subsection{Difference in Material Constitution}

Material difference is one important form of difference
through which identity continuity is often examined.

An existent may undergo material alteration through
the loss, replacement, addition, or reorganization
of material constituents.

Such changes may be extensive. Nevertheless, the
degree of material difference does not by itself
determine whether the identity of the existent has
persisted.

The preservation of material constituents does not
independently establish numerical identity. Different
existents may share the same material constituents
or material composition under a specified description.

Conversely, the replacement of some or many material
constituents does not automatically establish that
the existent has become numerically different.

The ontological significance of material transformation
depends upon the kind of existent, the relevant
conditions of its constitution, and the respect
under which identity is being examined.

Material difference must therefore be distinguished
from numerical non-identity.

\subsection{Difference in Structural Organization}

Structural difference concerns alterations in the
organization of components, relations, or internal
configurations.

An existent may undergo structural reorganization
while preserving some of its previous relations.
It may also acquire a substantially different
organization through transformation.

Neither structural preservation nor structural
difference independently determines numerical identity.

A similar structure may be realized by distinct
existents. A structure may also undergo transformation
without necessarily establishing the replacement
of the existent whose organization has changed.

The persistence of a structure and the persistence
of identity are therefore separate questions.

Structural difference may be relevant to an account
of identity when the structure has been shown to
play an ontologically significant role. Yet its
relevance cannot be assumed merely because structural
features are comparatively stable or readily described.

The transition from structural difference to
non-identity requires an account of why the particular
transformation is incompatible with the persistence
of the identity under examination.

\subsection{Difference in Functional Determination}

Functional difference concerns changes in the
operations, capacities, roles, or activities associated
with an existent.

An existent may acquire new functions, lose previous
functions, or undergo a transformation in the manner
in which its capacities are actualized.

Functional continuity does not independently establish
identity, since distinct existents may perform the
same function.

Likewise, functional discontinuity does not independently
establish non-identity. An existent may undergo
functional transformation without the numerical
identity question being settled by that transformation.

The same function may be realized through different
structures, materials, or relations. Conversely,
the same existent may participate in different
functional arrangements across successive conditions.

A functional description therefore cannot serve
as a universal criterion of numerical identity.

The ontological relevance of functional difference
must be established in relation to the existent
and the conditions under which its identity is
being examined.

\subsection{Difference in Relational Configuration}

An existent may undergo changes in its relational
configuration.

It may enter into new relations, lose previous
relations, or participate in a different organization
of relations while retaining other determinations.

Such relational transformations may significantly
alter its position within a wider configuration.
However, a change in relational position does not
automatically establish a change in numerical identity.

The same relational position may be occupied by
different existents. An existent may also occupy
different relational positions across successive
conditions.

Relational difference must therefore be distinguished
from numerical non-identity.

A relational configuration may contribute to the
conditions of identity preservation in a particular
domain, but its significance must be independently
established. It cannot be treated as a universal
criterion without examining the ontological status
of the relations involved.

\subsection{Difference in Manifestation}

Manifestations associated with an existent may
differ substantially across successive conditions.

An existent may manifest different determinations,
configurations, capacities, or relations. It may
also undergo changes in the extent, form, or organization
of its manifestations.

Such differences may establish that the manifestations
are not qualitatively identical. They do not
automatically establish that the identity associated
with them is numerically different.

The distinction between identity and manifestation
is therefore essential to the analysis of difference.

A manifestation cannot be treated as an exhaustive
representation of the identity of the existent.
Consequently, a difference between manifestations
cannot be transferred without argument into a
difference between identities.

Likewise, similar manifestations do not independently
establish that their sources or bearers possess
the same numerical identity.

The relation between manifestation difference and
identity difference must remain an independent
ontological question.

\subsection{The Difference Between Change and Replacement}

Change and replacement must not be treated as
equivalent.

Change concerns the alteration of a subject or
one of its determinations. Replacement concerns
the emergence or substitution of a numerically
different subject in relation to a previous one.

A change may be extensive without necessarily
constituting replacement. Conversely, replacement
may occur without requiring a substantial qualitative
difference between the previous and subsequent
subjects.

This distinction prevents the magnitude of change
from becoming an automatic criterion of non-identity.

The claim that an existent has changed significantly
does not, by itself, establish that another existent
has replaced it. Such a conclusion requires an
account of the ontological conditions under which
the transformation constitutes replacement rather
than alteration within the same identity.

Similarly, the restoration of a previous configuration
does not necessarily establish the restoration
of the previous numerical identity.

Change, replacement, and restoration must therefore
remain analytically distinct.

\subsection{The Problem of Gradual Difference}

Gradual transformation introduces a further difficulty.

An existent may undergo a sequence of relatively
small changes, each of which differs only slightly
from the preceding condition. Across an extended
succession, however, the final condition may differ
substantially from the initial one.

The gradual character of the transformation does
not independently establish continuity of identity.
Nor does the accumulation of qualitative differences
automatically establish numerical non-identity.

The problem cannot be resolved merely by examining
the magnitude of each individual change or the total
difference between the initial and final conditions.

A sequence of small transformations may involve
continuity of one existent, the succession of distinct
existents, or another ontological configuration
that requires further determination.

The question of identity must therefore be examined
across the succession as a whole and in relation
to the ontological conditions relevant to the
existent under consideration.

The absence of a single decisive transformation
does not establish identity continuity. Equally,
the accumulation of differences does not independently
establish identity discontinuity.

\subsection{Similarity and the Problem of False Identity}

The distinction between difference and non-identity
must be accompanied by a distinction between similarity
and identity.

Two existents may share a substantial number of
features, possess similar structures, perform
equivalent functions, or exhibit comparable
manifestations without being numerically identical.

Similarity concerns a relation of resemblance
under selected respects. Numerical identity concerns
whether the subjects are one and the same existent.

The two claims cannot be substituted for one another.

A sufficiently high degree of similarity does not
establish numerical identity. Nor does the reproduction
of a configuration or structure establish that
the reproduced existent is numerically identical
with the original.

This distinction is particularly important when
identity continuity is examined through preserved
features. A preserved feature may characterize
both successive conditions and distinct existents.

Consequently, the identification of similarity
or equivalence must not be treated as a complete
account of identity.

\subsection{Difference and the Problem of Multiple Identities}

Difference may also arise in relation to the possibility
of multiple identities.

An existent may exhibit several distinct configurations,
manifestations, or relational determinations. These
differences do not automatically establish that
the existent possesses multiple identities.

Conversely, the apparent unity of an existent does
not independently establish that only one identity
is involved.

The relation between difference and plurality of
identities must therefore be examined without
assuming that every distinct manifestation corresponds
to a distinct identity or that every manifestation
belongs to one identity.

The question of whether a transformation concerns
different determinations of one identity or the
relation between multiple identities cannot be
resolved through qualitative difference alone.

This issue requires an independent investigation
of the ontological conditions under which identity
is individuated and preserved.

The present section therefore leaves the number
of identities open rather than deriving it from
the multiplicity or similarity of observed conditions.

\subsection{Difference Without a Determinate Conclusion}

There are cases in which a difference can be established
without the available ontological grounds being
sufficient to determine whether identity has persisted.

A material, structural, functional, or relational
transformation may be clearly identifiable while
its significance for numerical identity remains
undetermined.

In such cases, the appropriate conclusion is not
necessarily identity or non-identity. It may be
that the relevant difference has been established
but that its ontological implications remain unresolved.

This distinction is methodologically important.

The inability to establish identity continuity
does not automatically establish non-identity.
Likewise, the inability to establish non-identity
does not automatically establish continuity.

An unresolved question must remain unresolved until
adequate grounds are available.

The distinction between established difference
and undetermined identity prevents the analysis
from producing conclusions that exceed the evidence
and ontological argument available.

\subsection{The Subject and Respect of Difference}

Every claim concerning difference must specify
its subject and respect.

The subject may be an existent, an identity,
a determination, a manifestation, a structure,
a relation, or a temporal condition.

The respect may concern material constitution,
structural organization, functional determination,
relational configuration, or another explicitly
defined domain.

A difference established between two manifestations
cannot automatically be transferred to the identities
associated with them. A difference in material
constitution cannot automatically be generalized
to every other respect of the existent.

Likewise, the comparison of two temporal conditions
must distinguish between differences within one
existent and differences between distinct existents.

Without these specifications, the concept of difference
remains too broad to support a determinate conclusion
about numerical identity.

\subsection{Methodological Requirements}

The investigation of identity and non-identity through
difference must observe the following requirements:

\begin{enumerate}
    \item Qualitative difference must be distinguished
    from numerical non-identity.
    \item Similarity must not be treated as sufficient
    proof of numerical identity.
    \item The magnitude of transformation must not
    be used as an automatic criterion of replacement.
    \item Material, structural, functional, relational,
    and manifestational differences must be examined
    separately.
    \item Change must be distinguished from replacement
    and restoration.
    \item Gradual transformation must not be assumed
    to guarantee identity continuity or discontinuity.
    \item Difference between manifestations must not
    be transferred automatically to difference between
    identities.
    \item The possibility of multiple identities must
    remain open to independent ontological investigation.
    \item Every conclusion must specify its subject,
    domain, and respect.
\end{enumerate}

These requirements prevent difference from becoming
an unrestricted criterion through which identity
is either affirmed or denied without adequate
ontological justification.

\subsection{Scope of the Present Section}

This section has established the distinction between
qualitative difference and numerical non-identity.

An existent may undergo substantial differences
in its material constitution, structural organization,
functional determination, relational configuration,
or manifestations without these differences alone
establishing that its identity has ceased.

Conversely, similarity between successive conditions
does not independently establish that they belong
to one and the same numerical identity.

Difference is therefore relevant to the investigation
of identity continuity, but it cannot serve as
an independent and sufficient criterion of either
identity preservation or non-identity.

The ontological significance of difference must be
determined in relation to the subject, domain, and
respect under examination.

The next section investigates the conditions of
identity preservation. It examines whether and under
what circumstances particular ontological conditions
can justify the claim that identity persists through
transformation, without reducing identity to any
single material, structural, relational, functional,
or invariant feature.

% ============================================================
% End of Section 04
% ============================================================